\documentclass[12pt]{amsart}
% Default margins are too wide all the way around. I reset them here
\setlength{\hoffset}{-.75in}
\setlength{\textwidth}{6.5in}
\setlength{\voffset}{-.5in}
\setlength{\textheight}{9.0in}
\usepackage{amsmath}
\usepackage{amssymb}
\usepackage{amsthm}
\usepackage{pgfplots}
\usepackage{enumerate}
\usepackage{tikz}
\usetikzlibrary{shadows,arrows,arrows.meta,shapes,positioning,calc}% Required for the Circle and Triangle arrow tips
\usepackage{setspace}
\usepackage{siunitx}
\usepackage{enumitem}
\usepackage{tcolorbox}
\usepackage{array}
\usepackage{microtype}

\theoremstyle{conjecture}
\newtheorem{conjecture}{Conjecture}

\theoremstyle{definition}
\newtheorem{definition}{Definition}

\theoremstyle{question}
\newtheorem{question}{Question}

\theoremstyle{theorem}
\newtheorem{theorem}{Theorem}

\newenvironment{solution}
{\begin{proof}[Solution]}
{\end{proof}}

\newcommand{\sectionrule}{%
    \vspace{0.3em}
    \hrule
    \vspace{0.8em}
}

\title{Dynamical Systems:\\Class August 27, 2026\\ 3.a Nondimensionalization.}
\author{Rob Tompkins}
\date{\today}


\begin{document}

    \maketitle

    \textbf{Extra vocabulary / extra facts:}

    \begin{tcolorbox}[
        colback=white,
        colframe=black!70,
        boxrule=0.8pt,
        arc=2mm,
        left=3mm,
        right=3mm,
        top=2mm,
        bottom=2mm
    ]
        \begin{itemize}
            \item A \textbf{nondimensional group} is a group of parameters or constants that together are
            dimensionless but that have the property that any factor of the group has dimension.

            \item A \textbf{Monod function} is a type of switching function (just as $\tanh x$ was an example of
            a switching function). The Monod function has the form
            \[
                f(x)=r\frac{x}{a+x}.
            \]

            \item A \textbf{Hill function} is a type of switching function (compare to $\tanh x$ and to the Monod
            function). The Hill function has the form
            \[
                f(x)=r\frac{x^n}{a^n+x^n}
            \]
            where $n$ is the \textit{Hill coefficient}.
        \end{itemize}
    \end{tcolorbox}

    \sectionrule

    {\Large\textbf{Skill check practice}}

    \begin{enumerate}
        \item Consider the differential equation
        \[
            \frac{dx}{d\tau}
            =
            rT_0\left(\frac{1}{h_v}+x\right)
            -
            \frac{rT_0A}{K}x.
        \]

        Assume that $x$ and $\tau$ are nondimensional variables and that $r,T_0,A,K$ are parameters with
        dimension.

        Identify two nondimensional groups from the equation above.
    \end{enumerate}

    \sectionrule

    {\Large\textbf{Skill check practice solution}}

    \begin{enumerate}
        \item $x$ and $\tau$ are nondimensional variables (that info is given). Recall that quantities that are added
        to each other or that are equal to each other must have the same dimensions.

        We have
        \[
            \left[\frac{dx}{d\tau}\right]
            =
            \left[
                rT_0\left(\frac{1}{h_v}+x\right)
            \right]
            =
            \left[
                \frac{rT_0A}{K}x
            \right]
        \]
        (where $[\cdot]$ denotes ``the dimensions of'').

        We have
        \[
            [x]=[\tau]=1.
        \]
        So
        \[
            \left[\frac{dx}{d\tau}\right]=1.
        \]
        That means every term in the equation is also all dimensionless.

        The dimension of a product is the product of the dimensions, so
        \[
            1
            =
            [rT_0]
            \left[
                \left(\frac{1}{h_v}+x\right)
            \right]
            =
            \left[
                \frac{rT_0A}{K}
            \right]
            [x].
        \]

        The dimension of a sum is the dimension of either component of the sum, so
        \[
            1
            =
            [rT_0][x]
            =
            \left[
                \frac{rT_0A}{K}
            \right]
            [x].
        \]

        \newpage

        Using $[x]=1$, this is
        \[
            1=[rT_0]
            =
            \left[
                \frac{rT_0A}{K}
            \right].
        \]

        $rT_0$ is a dimensionless group.
        \[
            1
            =
            \left[
                \frac{rT_0A}{K}
            \right]
            =
            [rT_0]
            \left[
                \frac{A}{K}
            \right].
        \]
        $\dfrac{A}{K}$ is another dimensionless group.

        $h_v$ is a dimensionless parameter (but is not a group).
    \end{enumerate}

    \sectionrule

    \newpage

    {\Large\textbf{Teams}}

    Post screenshots of your work to the course Google Drive today:Link to drive folder. Include words,
    labels, and other short notes on your whiteboard that might make those solutions useful to you or
    your classmates. Name your images using the format:

    \begin{itemize}
        \item \texttt{teamXX-YY.EXT}
    \end{itemize}

    where XX is your two-digit team number (e.g.\ 06 if you are team 6), YY is the two digit number of
    image for today (e.g.\ 01 for the first image, 02 for the second, and so on.), and EXT is the extension
    for the image filetype you used.

    \sectionrule

    {\LARGE\textbf{Team activity}}

    \vspace{0.7em}

    {\Large\textbf{Assigning roles}}

    \begin{itemize}
        \item Have a brief rock-paper-scissors tournament (best of 1).

        \item The winner will be the scribe for the first problem (doing all writing for the team).

        \item Re-match with the other team members. The winner will be the time-keeper (keeping the
        team on task and making sure that team members are taking turns contributed / sharing
        time).

        \item Choose a remaining team member to be the questioner (checking in to make sure that team
        members are asking any questions that they have).
    \end{itemize}

    \vspace{0.5em}

    \textbf{Nondimensionalization}

    \begin{enumerate}
        \item Let
        \[
            \dot N = rN(1-N/K)-H.
        \]
        This is a logistic population model where there is a constant harvesting rate reducing
        population growth.

        \begin{enumerate}
            \item For each of the variables and each of the parameters, identify its associated dimension.
            \textit{Write this out in expressions of the form $[a]=L$.}

            \item Create dimensional constants, $N_0$ and $T_0$, and use them to create nondimensional variables
            \[
                x=N/N_0
                \qquad\text{and}\qquad
                \tau=t/T_0.
            \]
            Substitute $x$ and $\tau$ into the equation and simplify.

            \item List all nondimensional groups that arise. Consider a combination a \textit{nondimensional group}
            if the combination is nondimensional but every piece would be dimensional if it were broken
            apart in any way.

            \item Make choices for values of the constants $T_0$ and $N_0$ that eliminate two of the nondimen-
            sional groups.

            \item Define a new nondimensional parameter (use a Greek letter such as $\alpha,\beta,\gamma,\mu$), and rewrite
            your equation as a nondimensional one.

            \item How many parameters are there in the nondimensional system? How does this compare
            to the number in the dimensional version? Notice that each dimensional constant we
            introduce enables us to remove a parameter, so that the nondimensional equation has
            fewer parameters than the dimensional one.
            \textit{This result on the reduction in the number of
            parameters from nondimensionalization is called the Buckingham Pi theorem. It is called}

            \newpage

            \textit{the ``pi'' theorem because he used the symbols $\Pi_1,\Pi_2$, etc to represent the nondimensional
            groups.}
        \end{enumerate}

        \item \textit{Rotate your roles: questioner $\rightarrow$ timekeeper $\rightarrow$ scribe}

        Let
        \[
            \dot N = rN(1-N/K)-HN/(A+N).
        \]
        This is a slightly different harvesting case.

        \begin{enumerate}
            \item This harvesting term, $HN/(A+N)$, is in the form of a special function called a
            \textbf{Monod function}. Plot an approximation of the Monod function by hand
            \textbf{without} using any plotting tools.

            \item What do you think of this function as a description of a harvesting process?

            \item Identify the dimension of each of the variables and parameters. Once you nondimensionalize
            how many parameters do you expect to remain?

            \item Nondimensionalize this equation.

            \begin{itemize}
                \item As part of this process, identify the dimensionless groups and check your work with
                another team.

                \item There are multiple good choices for $N_0$. What are some reasons to choose one or the
                other?
            \end{itemize}

            \item Now that it's nondimensional, take another look at the expression for your harvesting
            function. Plot it using appropriate axis ticks. How did your axis tick labels change?
        \end{enumerate}
    \end{enumerate}

    \textit{tips for plotting the monod function:}

    \begin{itemize}
        \item How does the function behave as $N\to0$?

        \item How does it behave as $N\to\infty$?

        \item Use a unit increment of $A$ on the horizontal axis and an increment of $H$ on the vertical
        axis. Mark the point corresponding to an input of $N=A$.

        \item Approximate the function for $N$ close to zero, using $A+N\approx A$ for $N$ small enough.

        \item Draw a curve that connects up the features you have identified.
    \end{itemize}

\end{document}